% miniVivisection -- title page designs for the manual (LaTeX)
%
% Created:   2026-10-19T19:41:44+03:00
% Updated:   2026-10-19T20:21:14+03:00 -- designs 17-25, the outlined V° in the typographic layout
% Generator: Claude Code (Model: Claude Opus 5.5, claude-opus-5-5)
%
% Original prompt:
%   in a seperate latex /pdf file generate a varaiety of different designs for the first page, the
%   title page, of the manual, with color, whiou out color, with whithout screen shot, with the V° logo
%   as the main visual elemtn, etc.
%
% Context:
% - Twenty-five alternative title pages for miniVivisection_Manual.tex, one per page, with an overview:
%   color and black-and-white, with and without screenshots, the V° mark as the main element, and
%   themes drawn from the plug-in itself (name dots, a Rhino viewport, the code editor, a blueprint).
% - Build: python tools/make_manual_latex.py --titles -> docs/manual/title_designs.pdf. To use one,
%   copy its picture into chapters/00_title.tex.

\documentclass[a4paper, 10.5pt, oneside, open=any, parskip=false]{scrbook}
% miniVivisection manual -- the look: fonts, colors, headings, boxes, keys, code.
% Compiled with LuaLaTeX (fontspec, system fonts for symbols): tools/make_manual_latex.py.

% ---------------------------------------------------------------- fonts
\usepackage{fontspec}
\setmainfont{SourceSerif4}[
  UprightFont = *-Regular, ItalicFont = *-RegularIt, BoldFont = *-Semibold, BoldItalicFont = *-SemiboldIt,
  Extension = .otf, Numbers = Lining]
\setsansfont{SourceSans3}[
  UprightFont = *-Regular, ItalicFont = *-RegularIt, BoldFont = *-Semibold, BoldItalicFont = *-SemiboldIt,
  Extension = .otf]
\setmonofont{SourceCodePro}[
  UprightFont = *-Regular, ItalicFont = *-RegularIt, BoldFont = *-Semibold,
  Extension = .otf, Scale = 0.86]
\newfontfamily\mvhead{SourceSans3}[
  UprightFont = *-Semibold, BoldFont = *-Bold, ItalicFont = *-SemiboldIt, Extension = .otf, Ligatures = TeX]
\newfontfamily\mvlight{SourceSans3}[UprightFont = *-Light, Extension = .otf, Ligatures = TeX]
\newfontfamily\symbolfont{Segoe UI Symbol}
% The V° mark, in the font of the V° button on Grasshopper's toolbar.
\newfontfamily\vlogo{Segoe UI Bold}
% Outlined text: stroked, not filled (PDF text render mode 1), the stroke #1 bp wide.
\providecommand{\outlined}[2]{\pdfextension literal{1 Tr #1 w}#2\pdfextension literal{0 Tr}}
\usepackage{newunicodechar}
\newunicodechar{◆}{{\symbolfont ◆}}
\newunicodechar{☐}{{\symbolfont ☐}}
\newunicodechar{☑}{{\symbolfont ☑}}
\newunicodechar{∣}{{\symbolfont ∣}}
\newunicodechar{✓}{{\symbolfont ✓}}
\newunicodechar{▸}{{\symbolfont ▸}}
\newunicodechar{⟶}{{\symbolfont ⟶}}
\usepackage{microtype}
\usepackage[strings]{underscore} % _ in text: code names such as gh_item

% ---------------------------------------------------------------- colors
\usepackage[table]{xcolor}
\definecolor{mvblue}{HTML}{1A4D8F}
\definecolor{mvbluelight}{HTML}{EAF0F8}
\definecolor{mvbluemid}{HTML}{C9D6EA}
\definecolor{mvorange}{HTML}{D9661F}
\definecolor{mvorangelight}{HTML}{FCEFE6}
\definecolor{mvgreen}{HTML}{2E7D4F}
\definecolor{mvgreenlight}{HTML}{E8F4EC}
\definecolor{mvred}{HTML}{B3261E}
\definecolor{mvredlight}{HTML}{FBEAE9}
\definecolor{mvgray}{HTML}{5F6B7A}
\definecolor{mvgraylight}{HTML}{F2F4F7}
\definecolor{mvink}{HTML}{1D1D1F}
\definecolor{codebg}{HTML}{F6F8FA}
\definecolor{codekw}{HTML}{0000FF}
\definecolor{codestr}{HTML}{A31515}
\definecolor{codecom}{HTML}{008000}
\definecolor{codenum}{HTML}{098658}
\definecolor{codefn}{HTML}{795E26}
\color{mvink}

% ---------------------------------------------------------------- page
\usepackage[a4paper, top=24mm, bottom=26mm, inner=24mm, outer=24mm, headsep=9mm, footskip=12mm]{geometry}
\setlength{\parindent}{0pt}
\setlength{\parskip}{0.55em plus 0.2em minus 0.1em}
\linespread{1.07}
\usepackage{scrlayer-scrpage}
\pagestyle{scrheadings}
\clearpairofpagestyles
\automark[section]{chapter}
\ihead{\sffamily\small\color{mvgray}miniVivisection \,·\, Manual}
\ohead{\sffamily\small\color{mvgray}\leftmark}
\ofoot*{\sffamily\small\color{mvblue}\thepage}
\setkomafont{pageheadfoot}{\normalfont}
\renewcommand*{\chaptermarkformat}{\thechapter\enskip}

% ---------------------------------------------------------------- headings
\setkomafont{disposition}{\mvhead\color{mvblue}}
\setkomafont{chapter}{\mvhead\Huge\color{mvblue}}
\setkomafont{section}{\mvhead\Large\color{mvblue}}
\setkomafont{subsection}{\mvhead\large\color{mvink}}
\setkomafont{subsubsection}{\mvhead\normalsize\color{mvink}}
\setkomafont{paragraph}{\mvhead\normalsize\color{mvblue}}
\setkomafont{part}{\mvhead\Huge\color{mvblue}}
\setkomafont{partnumber}{\mvlight\Large\color{mvgray}}
\RedeclareSectionCommand[beforeskip=-1.6em plus -0.4em, afterskip=0.6em]{section}
\RedeclareSectionCommand[beforeskip=-1.2em plus -0.3em, afterskip=0.35em]{subsection}
\RedeclareSectionCommand[beforeskip=-0.9em plus -0.2em, afterskip=0.25em]{subsubsection}
\RedeclareSectionCommand[beforeskip=0.7em plus 0.2em, afterskip=-0.6em]{paragraph}
\RedeclareSectionCommand[beforeskip=0pt, afterskip=2.2em]{chapter}
\renewcommand*{\chapterformat}{%
  \begin{tikzpicture}[baseline=(n.base)]
    \node[fill=mvblue, text=white, inner xsep=7pt, inner ysep=5pt, rounded corners=3pt, font=\mvhead\huge] (n) {\thechapter};
  \end{tikzpicture}\hspace{0.6em}}
\renewcommand*{\chapterlinesformat}[3]{%
  \parbox{\linewidth}{\raggedright#2#3\par\vspace{0.35em}{\color{mvbluemid}\rule{\linewidth}{1.2pt}}}}
\renewcommand*{\sectionformat}{{\color{mvbluemid}\thesection}\hspace{0.7em}}

% ---------------------------------------------------------------- lists, tables, figures
\usepackage{enumitem}
\setlist{itemsep=0.2em, topsep=0.3em, parsep=0.15em, leftmargin=1.4em}
\setlist[itemize,1]{label={\color{mvblue}\small\textbullet}}
\setlist[itemize,2]{label={\color{mvgray}\footnotesize–}}
\setlist[description]{font=\sffamily\bfseries\color{mvink}, itemsep=0.25em, leftmargin=1.4em}
\usepackage{booktabs, tabularx, longtable, xltabular, array, multirow}
\renewcommand{\arraystretch}{1.22}
\newcolumntype{L}{>{\raggedright\arraybackslash}X}
\newcolumntype{P}[1]{>{\raggedright\arraybackslash}p{#1}}
\setlength{\aboverulesep}{0.25ex}
\setlength{\belowrulesep}{0.45ex}
\arrayrulecolor{mvbluemid}
\usepackage{graphicx}
% The screenshots as high-quality JPEGs, made from docs/screenshots/*.png by tools/make_manual_latex.py.
\graphicspath{{generated/jpg/}}
\usepackage[font={small,sf}, labelfont={bf,color=mvblue}, justification=raggedright, singlelinecheck=false, skip=6pt]{caption}
\usepackage{float}

% ---------------------------------------------------------------- tikz, boxes
\usepackage{tikz}
\usetikzlibrary{arrows.meta, positioning, shapes.geometric, calc, fit, backgrounds}
% TikZ's ragged alignment sets \spaceskip for Computer Modern -- far too wide in Source Sans: the font's own spaces.
\tikzset{every node/.append style={execute at begin node={\spaceskip=0pt\relax\xspaceskip=0pt\relax}}}
\usepackage[most]{tcolorbox}
\tcbset{
  mvbox/.style={enhanced, breakable, boxrule=0pt, arc=2pt, left=10pt, right=10pt, top=6pt, bottom=6pt,
    before skip=0.8em, after skip=0.8em, fonttitle=\sffamily\bfseries\small, coltitle=#1, attach title to upper={\enskip}},
}
\newtcolorbox{note}[1][Note]{mvbox=mvblue, colback=mvbluelight, borderline west={3pt}{0pt}{mvblue}, title={#1}}
\newtcolorbox{tip}[1][Tip]{mvbox=mvgreen, colback=mvgreenlight, borderline west={3pt}{0pt}{mvgreen}, title={#1}}
\newtcolorbox{caution}[1][Careful]{mvbox=mvorange, colback=mvorangelight, borderline west={3pt}{0pt}{mvorange}, title={#1}}
\newtcolorbox{lesson}[1]{mvbox=mvgray, colback=mvgraylight, borderline west={3pt}{0pt}{mvgray}, title={Try it: #1}}
% A concept: a titled box of its own.
\newtcolorbox{concept}[1]{enhanced, breakable, colback=white, colframe=mvbluemid, boxrule=0.8pt, arc=3pt,
  left=10pt, right=10pt, top=8pt, bottom=8pt, before skip=1em, after skip=1em,
  title={#1}, fonttitle=\mvhead\normalsize, coltitle=mvblue, colbacktitle=mvbluelight,
  attach boxed title to top left={xshift=8pt, yshift=-8pt}, boxed title style={boxrule=0.8pt, colframe=mvbluemid, arc=2pt}}

% ---------------------------------------------------------------- inline markup
% \vn{Watch} -> V°Watch, the plug-in's V names.
\newcommand{\vn}[1]{\mbox{\sffamily\bfseries V°#1}}
% \ui{Edit in Rhino}: a label, button or menu item as the user sees it.
\newcommand{\ui}[1]{{\sffamily\bfseries #1}}
% \menu{A}{B}: a path through menus or settings groups.
\newcommand{\menu}[1]{{\sffamily\bfseries\def\>{\,{\color{mvgray}\small▸}\,}#1}}
% \setting{V°Settings > Debugger > Stop once per line}
\newcommand{\setting}[1]{{\sffamily\itshape\color{mvblue}#1}}
% \code{...}: code in the text.
\newcommand{\code}[1]{\texttt{\color{mvink}#1}}
\setlength{\emergencystretch}{2.5em}
\newcommand{\codebreak}[1]{\texttt{\color{mvink}#1}}
% \bs: a backslash in a path, where the line may break.
\newcommand{\bs}{\textbackslash\allowbreak}
% \file{...}: file and folder names.
\newcommand{\file}[1]{\texttt{\color{mvgray}#1}}
% \key{Ctrl}: a key cap; \keys{Ctrl+Alt+B} for a chord.
\newcommand{\key}[1]{\tikz[baseline=(k.base)]\node[draw=mvgray!70, fill=mvgraylight, rounded corners=1.5pt, inner xsep=2.6pt, inner ysep=1.4pt, line width=0.45pt, font=\sffamily\scriptsize] (k) {\strut #1};}
\ExplSyntaxOn
\NewDocumentCommand{\keys}{m}{
  \seq_set_split:Nnn \l_tmpa_seq {+} {#1}
  \bool_set_true:N \l_tmpa_bool
  \seq_map_inline:Nn \l_tmpa_seq {
    \bool_if:NF \l_tmpa_bool { \mbox{\scriptsize\,+\,} }
    \bool_set_false:N \l_tmpa_bool
    \key{##1}
  }
}
\ExplSyntaxOff
% A reference to a section, shown as "§ 5.2" with a link.
\newcommand{\see}[1]{\mbox{\hyperref[#1]{§\,\ref*{#1}}}}
% A lesson's file.
\newcommand{\lessonfile}[1]{\texttt{\color{mvgreen}#1}}

% ---------------------------------------------------------------- code listings
\usepackage{listings}
\lstdefinestyle{mv}{
  basicstyle=\ttfamily\small, columns=fullflexible, keepspaces=true,
  backgroundcolor=\color{codebg}, frame=single, framerule=0pt, rulecolor=\color{codebg},
  framesep=6pt, xleftmargin=6pt, xrightmargin=6pt, aboveskip=0.8em, belowskip=0.8em,
  keywordstyle=\color{codekw}, stringstyle=\color{codestr}, commentstyle=\color{codecom}\itshape,
  showstringspaces=false, upquote=true, breaklines=true, breakatwhitespace=true,
  literate={°}{{°}}1 {→}{{$\rightarrow$}}1 {…}{{\dots}}1 {◆}{{\symbolfont ◆}}1,
}
\lstdefinelanguage{PythonMV}[]{Python}{morekeywords={True,False,None,as,with,async,await,nonlocal}}
\lstset{style=mv, language=PythonMV}
\lstdefinestyle{plain}{style=mv, language={}}

% ---------------------------------------------------------------- links
\usepackage[hidelinks, bookmarksnumbered, bookmarksopen, bookmarksopenlevel=1, pdfdisplaydoctitle]{hyperref}
\usepackage{bookmark}
\hypersetup{colorlinks=true, linkcolor=mvblue, urlcolor=mvblue, citecolor=mvblue}
\usepackage[capitalise, noabbrev]{cleveref}

% ---------------------------------------------------------------- table of contents
\setkomafont{chapterentry}{\mvhead\color{mvblue}}
\setkomafont{partentry}{\mvhead\large\color{mvink}}
\RedeclareSectionCommand[tocentryformat=\mvhead\color{mvblue}, tocpagenumberformat=\mvhead\color{mvblue}]{chapter}
\setcounter{tocdepth}{1}

% A figure with one of the screenshots (960 x 600): \shot{file}{caption}{label}
\newcommand{\shot}[3]{%
  \begin{figure}[htbp]
    \centering
    \tikz\node[inner sep=0pt, draw=mvbluemid, line width=0.6pt]{\includegraphics[width=0.92\linewidth]{#1}};
    \caption{#2}\label{#3}
  \end{figure}}

% A feature list's table: \begin{features}{Heading} \feat{name}{what}{section} ... \end{features}
\usepackage{needspace}
\newenvironment{features}[1]{%
  \par\needspace{7\baselineskip}\smallskip{\mvhead\color{mvblue}#1}\par\nopagebreak\vspace{-0.4em}
  \xltabular{\linewidth}{@{}P{0.27\linewidth}L>{\raggedleft\arraybackslash}p{0.07\linewidth}@{}}
  \toprule}{\bottomrule\endxltabular}
\newcommand{\feat}[3]{{\sffamily\bfseries #1} & #2 & \see{#3}\\}

\input{generated/version}
\usetikzlibrary{fadings, shadows}
\hypersetup{pdftitle={miniVivisection -- title page designs}}

% \vlogo (the V° mark's font) and \outlined come from style.tex.

% The texts every design uses.
\newcommand{\mvTitle}{miniVivisection}
\newcommand{\mvSub}{See your script's data in Rhino while you debug}
\newcommand{\mvWhat}{A plug-in for the Script Editor and Grasshopper of Rhino~8}
\newcommand{\mvVer}{Manual \,·\, Version \mvversion\ (build \mvbuild) \,·\, \mvdate}

% One design: \design{name}{tikz drawing, in mm from the page's lower left corner}
\newcounter{design}
\newcommand{\design}[2]{%
  \clearpage\thispagestyle{empty}\stepcounter{design}\null
  \begin{tikzpicture}[remember picture, overlay, x=1mm, y=1mm]
    \begin{scope}[shift={(current page.south west)}]
      #2
      \node[anchor=south, inner sep=1pt, fill=white, fill opacity=0.85, text opacity=1]
        at (105,2.5) {\sffamily\tiny\color{mvgray}Design \thedesign\ · #1};
    \end{scope}
  \end{tikzpicture}}

% Viewport colors, as Rhino draws its Top view, and the watch color of the screenshots.
\definecolor{rhinobg}{HTML}{C8C8C8}
\definecolor{rhinominor}{HTML}{B4B4B4}
\definecolor{rhinomajor}{HTML}{9A9A9A}
\definecolor{watchmag}{HTML}{D63AD6}
% VS Code's Dark+ colors, for the code editor design.
\definecolor{vsbg}{HTML}{1E1E1E}
\definecolor{vsgutter}{HTML}{858585}
\definecolor{vskw}{HTML}{C586C0}
\definecolor{vsblue}{HTML}{569CD6}
\definecolor{vsstr}{HTML}{CE9178}
\definecolor{vscom}{HTML}{6A9955}
\definecolor{vsfn}{HTML}{DCDCAA}
\definecolor{vsvar}{HTML}{9CDCFE}
\definecolor{vsnum}{HTML}{B5CEA8}
\definecolor{vstext}{HTML}{D4D4D4}
\definecolor{vsstatus}{HTML}{007ACC}
% The engraving plate.
\definecolor{paper}{HTML}{F4EDDC}
\definecolor{sepia}{HTML}{3B2A1A}
% The blueprint.
\definecolor{blueprint}{HTML}{0E3A6B}

\begin{document}

% ------------------------------------------------------------------ the overview
\thispagestyle{empty}
{\mvhead\Huge\color{mvblue} Title page designs}\par\vspace{0.4em}
{\sffamily\color{mvgray} Twenty-five alternatives for the first page of the miniVivisection manual, one per page.
To use one, copy its drawing into \file{chapters/00_title.tex}.}\par\vspace{1.2em}
\newcommand{\yes}{\textcolor{mvgreen}{✓}}
{\small
\begin{xltabular}{\linewidth}{@{}>{\sffamily}r>{\sffamily\bfseries}P{0.29\linewidth}>{\centering\arraybackslash}p{0.06\linewidth}>{\centering\arraybackslash}p{0.06\linewidth}>{\centering\arraybackslash}p{0.06\linewidth}L@{}}
\toprule
 & \textsf{\textbf{Design}} & \textsf{\textbf{Color}} & \textsf{\textbf{Shot}} & \textsf{\textbf{V°}} & \textsf{\textbf{Idea}} \\ \midrule
1  & Blue band and screenshot    & \yes & \yes &      & the current title page \\
2  & Black band, grayscale       &      & \yes &      & design 1 without color, for printing \\
3  & Typographic, black and white &     &      &      & Swiss: type, rules and white space only \\
4  & The V° mark, large          & \yes &      & \yes & the mark as the cover's picture \\
5  & The V° mark on blue         & \yes &      & \yes & reversed, the mark bleeding off the page \\
6  & The V° mark in outline      &      &      & \yes & line only, black and white \\
7  & The toolbar button          & \yes &      & \yes & the V° button on Grasshopper's toolbar, enlarged \\
8  & Color strip with the mark   & \yes & \yes & \yes & a blue strip on the left, a screenshot on the right \\
9  & Screenshot, full bleed      & \yes & \yes &      & the picture first, the title on a dark band \\
10 & Mosaic of screenshots       & \yes & \yes &      & all sixteen screenshots as a wall \\
11 & Mosaic, grayscale           &      & \yes &      & the wall without color \\
12 & Name dot                    & \yes &      &      & the title in a name dot, pointing at a curve \\
13 & Rhino viewport              & \yes &      &      & lesson 01's bow tie, drawn in a Top view \\
14 & Code editor                 & \yes &      &      & the title as a paused script, dark theme \\
15 & Blueprint                   & \yes &      & \yes & the mark as a technical drawing, with a title block \\
16 & Anatomical plate            & \yes &      & \yes & a vivisection: the mark laid open, like an engraving \\ \midrule
\multicolumn{6}{@{}l}{\textsf{\textbf{Fusions: the outlined V° (6) in the typographic layout (3)}}} \\ \midrule
17 & Outlined V°, lower right    &      &      & \yes & the mark between the title and the footer \\
18 & Behind the type, bleeding   &      &      & \yes & a huge light outline under the type \\
19 & Blue V, orange °            & \yes &      & \yes & the outline in the manual's two colors \\
20 & V° on top, type below       &      &      & \yes & the order turned round \\
21 & Filled light gray, outlined &      &      & \yes & a gray body under the line \\
22 & On the grid                 &      &      & \yes & the layout's columns and baselines shown \\
23 & Reversed                    &      &      & \yes & white on black \\
24 & Only the ° in orange        & \yes &      & \yes & one color, one spot \\
25 & Inline lettering            &      &      & \yes & a white line inside a black one \\
\bottomrule
\end{xltabular}}
{\sffamily\footnotesize\color{mvgray} Shot: with a screenshot. V°: the V° mark as the main element, in the font of the V° button
on Grasshopper's toolbar. The grayscale designs (2, 11) turn the screenshots gray with a PDF blend mode; printed or in any
current PDF viewer they show without color.}

% ------------------------------------------------------------------ 1 blue band and screenshot
\design{Blue band and screenshot (the current title page)}{
  \fill[mvblue] (0,189) rectangle (210,297);
  \fill[mvorange] (0,186.5) rectangle (210,189);
  \node[anchor=north west, text=white, inner sep=0pt] at (24,267) {\mvhead\fontsize{42}{46}\selectfont\mvTitle};
  \node[anchor=north west, text=mvbluemid, inner sep=0pt] at (24.5,247) {\mvlight\fontsize{17}{22}\selectfont\mvSub};
  \node[anchor=north west, text=white, inner sep=0pt] at (24.5,223) {\sffamily\large\mvWhat};
  \node[anchor=north west, text=mvbluemid, inner sep=0pt] at (24.5,209) {\sffamily\mvVer};
  \node[anchor=north, inner sep=0pt, draw=mvbluemid, line width=0.8pt] at (105,171) {\includegraphics[width=150mm]{01_watch_values_in_rhino}};
  \node[anchor=north west, text width=150mm, inner sep=0pt] at (30,72) {\sffamily\small\color{mvgray}A paused script's
    \code{outline} and \code{corners} drawn in Rhino's Top view, each with its name: the outline crosses itself --
    two corners are swapped.};
}

% ------------------------------------------------------------------ 2 black band, grayscale
\design{Black band, grayscale screenshot (design 1 without color)}{
  \fill[black] (0,189) rectangle (210,297);
  \fill[black!45] (0,186.5) rectangle (210,189);
  \node[anchor=north west, text=white, inner sep=0pt] at (24,267) {\mvhead\fontsize{42}{46}\selectfont\mvTitle};
  \node[anchor=north west, text=black!25, inner sep=0pt] at (24.5,247) {\mvlight\fontsize{17}{22}\selectfont\mvSub};
  \node[anchor=north west, text=white, inner sep=0pt] at (24.5,223) {\sffamily\large\mvWhat};
  \node[anchor=north west, text=black!25, inner sep=0pt] at (24.5,209) {\sffamily\mvVer};
  \node[anchor=north, inner sep=0pt] (shot) at (105,171) {\includegraphics[width=150mm]{01_watch_values_in_rhino}};
  \fill[white, blend mode=saturation] (shot.south west) rectangle (shot.north east);
  \draw[black!40, line width=0.6pt] (shot.south west) rectangle (shot.north east);
  \node[anchor=north west, text width=150mm, inner sep=0pt] at (30,72) {\sffamily\small\color{black!60}A paused script's
    \code{outline} and \code{corners} drawn in Rhino's Top view, each with its name: the outline crosses itself --
    two corners are swapped.};
}

% ------------------------------------------------------------------ 3 typographic
\design{Typographic, black and white}{
  \node[anchor=south west, inner sep=0pt] at (24,270) {\sffamily\small\addfontfeature{LetterSpace=18}\bfseries MANUAL};
  \node[anchor=south east, inner sep=0pt] at (186,270) {\sffamily\small Version \mvversion};
  \draw[black, line width=0.5pt] (24,266) -- (186,266);
  \node[anchor=north west, inner sep=0pt] at (23,252) {\fontsize{58}{60}\selectfont{\mvlight mini}{\mvhead Vivisection}};
  \node[anchor=north west, inner sep=0pt, text width=162mm, align=left] at (24.5,222) {\rmfamily\itshape\fontsize{20}{26}\selectfont See your script's data in Rhino\\while you debug};
  \draw[black, line width=2.5pt] (24.5,196) -- (54.5,196);
  \node[anchor=north west, inner sep=0pt, text width=100mm] at (24.5,188) {\rmfamily\large A plug-in for the Script Editor\\
    and Grasshopper of Rhino~8};
  \draw[black, line width=0.5pt] (24,58) -- (186,58);
  \node[anchor=north west, inner sep=0pt, text width=50mm] at (24,53) {\sffamily\footnotesize\bfseries Rhino 8\\\mdseries Script Editor\\Grasshopper};
  \node[anchor=north west, inner sep=0pt, text width=50mm] at (80,53) {\sffamily\footnotesize\bfseries Version \mvversion\\\mdseries build \mvbuild\\experimental};
  \node[anchor=north west, inner sep=0pt, text width=50mm] at (136,53) {\sffamily\footnotesize\bfseries \mvdate\\\mdseries formerly\\ScriptEditorWatchViz};
}

% ------------------------------------------------------------------ 4 the V° mark, large
\design{The V° mark, large, in color}{
  \node[inner sep=0pt] at (105,196) {\vlogo\fontsize{330}{330}\selectfont{\color{mvblue}V}{\color{mvorange}°}};
  \node[anchor=north, inner sep=0pt] at (105,104) {\mvhead\fontsize{40}{44}\selectfont\color{mvink}\mvTitle};
  \node[anchor=north, inner sep=0pt] at (105,86) {\mvlight\fontsize{16}{20}\selectfont\color{mvgray}\mvSub};
  \draw[mvorange, line width=1.2pt] (95,72) -- (115,72);
  \node[anchor=north, inner sep=0pt] at (105,64) {\sffamily\color{mvink}\mvWhat};
  \node[anchor=south, inner sep=0pt] at (105,20) {\sffamily\small\color{mvgray}\mvVer};
}

% ------------------------------------------------------------------ 5 the V° mark on blue
\design{The V° mark on blue, bleeding off the page}{
  \fill[mvblue] (0,0) rectangle (210,297);
  \begin{scope}
    \clip (0,0) rectangle (210,297);
    \node[anchor=south east, inner sep=0pt, text=white, text opacity=0.13] at (232,-24)
      {\vlogo\fontsize{560}{560}\selectfont V°};
  \end{scope}
  \node[anchor=north west, inner sep=0pt] at (24,270) {\vlogo\fontsize{40}{40}\selectfont\color{white}V{\color{mvorange}°}};
  \node[anchor=north west, inner sep=0pt, text=white] at (24,236) {\mvhead\fontsize{46}{50}\selectfont\mvTitle};
  \node[anchor=north west, inner sep=0pt, text=mvbluemid, text width=140mm] at (24.5,214) {\mvlight\fontsize{18}{23}\selectfont\mvSub};
  \fill[mvorange] (24.5,190) rectangle (44.5,191.5);
  \node[anchor=north west, inner sep=0pt, text=white] at (24.5,184) {\sffamily\large\mvWhat};
  \node[anchor=south west, inner sep=0pt, text=mvbluemid] at (24.5,22) {\sffamily\mvVer};
}

% ------------------------------------------------------------------ 6 the V° mark in outline
\design{The V° mark in outline, black and white}{
  \node[inner sep=0pt, text=black] at (105,190) {\vlogo\fontsize{330}{330}\selectfont\outlined{1.4}{V°}};
  \draw[black, line width=0.4pt] (40,108) -- (170,108);
  \node[anchor=north, inner sep=0pt] at (105,100) {\mvhead\fontsize{38}{42}\selectfont\color{black}\mvTitle};
  \node[anchor=north, inner sep=0pt] at (105,82) {\rmfamily\itshape\fontsize{16}{20}\selectfont\color{black!70}\mvSub};
  \draw[black, line width=0.4pt] (40,70) -- (170,70);
  \node[anchor=north, inner sep=0pt] at (105,62) {\sffamily\color{black}\mvWhat};
  \node[anchor=south, inner sep=0pt] at (105,20) {\sffamily\small\color{black!60}\mvVer};
}

% ------------------------------------------------------------------ 7 the toolbar button
\design{The toolbar button, enlarged}{
  \fill[mvgraylight] (0,0) rectangle (210,297);
  % the button: a rounded square as Grasshopper's toolbar draws it, the V°, the green Run badge
  \fill[black, opacity=0.10, rounded corners=9mm] (68,155) rectangle (148,235);
  \filldraw[fill=white, draw=black!25, line width=0.8pt, rounded corners=9mm] (65,158) rectangle (145,238);
  \node[inner sep=0pt] at (103,200) {\vlogo\fontsize{150}{150}\selectfont\color{black!78}V°};
  \fill[mvgreen] (135,168) circle (14);
  \fill[white] (130.5,160.5) -- (130.5,175.5) -- (143,168) -- cycle;
  \node[anchor=north, inner sep=0pt] at (105,126) {\mvhead\fontsize{40}{44}\selectfont\color{mvink}\mvTitle};
  \node[anchor=north, inner sep=0pt] at (105,108) {\mvlight\fontsize{16}{20}\selectfont\color{mvgray}\mvSub};
  % the three modes
  \foreach \x/\label/\col in {65/Normal/black!30, 105/Stop and record/mvgreen, 145/Always debug/mvorange} {
    \filldraw[fill=white, draw=black!20, rounded corners=1.5mm] (\x-7,74) rectangle (\x+7,88);
    \node[inner sep=0pt] at (\x-1,81) {\vlogo\fontsize{18}{18}\selectfont\color{black!70}V°};
    \fill[\col] (\x+4.5,76.5) circle (2.4);
    \node[anchor=north, inner sep=0pt] at (\x,71) {\sffamily\scriptsize\color{mvgray}\label};
  }
  \node[anchor=south, inner sep=0pt] at (105,22) {\sffamily\small\color{mvgray}\mvVer};
}

% ------------------------------------------------------------------ 8 color strip with the mark
\design{Color strip with the mark, and a screenshot}{
  \fill[mvblue] (0,0) rectangle (70,297);
  \fill[mvorange] (70,0) rectangle (72,297);
  \node[anchor=north, inner sep=0pt] at (35,272) {\vlogo\fontsize{110}{110}\selectfont\color{white}V{\color{mvorange}°}};
  \node[rotate=90, anchor=west, inner sep=0pt] at (37,40) {\sffamily\addfontfeature{LetterSpace=25}\bfseries\color{mvbluemid}MANUAL};
  \node[anchor=south, inner sep=0pt, text=mvbluemid, align=center] at (35,18) {\sffamily\scriptsize Version \mvversion\\build \mvbuild\\\mvdate};
  \node[anchor=north west, inner sep=0pt] at (84,262) {\mvhead\fontsize{34}{38}\selectfont\color{mvblue}\mvTitle};
  \node[anchor=north west, inner sep=0pt, text width=112mm] at (84.5,246) {\mvlight\fontsize{15}{19}\selectfont\color{mvink}\mvSub};
  \node[anchor=north west, inner sep=0pt] at (84.5,226) {\sffamily\color{mvgray}\mvWhat};
  \node[anchor=north west, inner sep=0pt, draw=mvbluemid, line width=0.6pt] (s) at (84,206) {\includegraphics[width=112mm]{02_values_in_the_code}};
  \node[anchor=north west, inner sep=0pt, draw=mvbluemid, line width=0.6pt] at (84,132) {\includegraphics[width=112mm]{16_grasshopper_prompts}};
  \node[anchor=north west, inner sep=0pt, text width=112mm] at (84,59) {\sffamily\scriptsize\color{mvgray}Above: values in the
    code of a paused script, and drawn in Rhino. Below: a Grasshopper definition built by V°Claude from one message.};
}

% ------------------------------------------------------------------ 9 screenshot, full bleed
\design{Screenshot, full bleed, the title on a dark band}{
  \begin{scope}
    \clip (0,96) rectangle (210,297);
    \node[anchor=west, inner sep=0pt] at (-6,196.5) {\includegraphics[height=201mm]{01_watch_values_in_rhino}};
  \end{scope}
  \fill[mvink] (0,0) rectangle (210,96);
  \fill[mvink, path fading=north] (0,96) rectangle (210,150);
  \fill[mvorange] (24,80) rectangle (44,81.5);
  \node[anchor=north west, inner sep=0pt, text=white] at (24,74) {\mvhead\fontsize{44}{48}\selectfont\mvTitle};
  \node[anchor=north west, inner sep=0pt, text=white!75!mvink] at (24.5,54) {\mvlight\fontsize{17}{22}\selectfont\mvSub};
  \node[anchor=north west, inner sep=0pt, text=white!60!mvink] at (24.5,38) {\sffamily\mvWhat};
  \node[anchor=south west, inner sep=0pt, text=white!45!mvink] at (24.5,14) {\sffamily\small\mvVer};
}

% ------------------------------------------------------------------ 10, 11 mosaics
\newcommand{\mosaic}{%
  \foreach \name [count=\i from 0] in {01_watch_values_in_rhino, 05_which_branch_failed, 09_markdown_panel_links,
      13_markdown_text_link_hint, 02_values_in_the_code, 06_paused_run_badge, 10_link_followed_in_grasshopper,
      14_markdown_text_dialog, 03_python_in_the_paused_frame, 07_script_panel_as_code, 11_v_edit, 15_link_followed_in_rhino,
      04_stop_and_record, 08_panel_edited_at_the_stop, 12_v_text, 16_grasshopper_prompts, 03_python_in_the_paused_frame,
      09_markdown_panel_links, 11_v_edit, 01_watch_values_in_rhino, 06_paused_run_badge} {
    \pgfmathtruncatemacro{\r}{\i/3}
    \pgfmathtruncatemacro{\c}{mod(\i,3)}
    \node[anchor=north west, inner sep=0pt] at (\c*70.5-0.5, 297.5-\r*44.25) {\includegraphics[width=69.5mm]{\name}};
  }}
\design{Mosaic of the screenshots}{
  \fill[white] (0,0) rectangle (210,297);
  \mosaic
  \fill[mvblue, opacity=0.93] (0,104) rectangle (210,196);
  \fill[mvorange] (0,102) rectangle (210,104);
  \node[anchor=north west, inner sep=0pt, text=white] at (24,182) {\mvhead\fontsize{42}{46}\selectfont\mvTitle};
  \node[anchor=north west, inner sep=0pt, text=mvbluemid] at (24.5,162) {\mvlight\fontsize{17}{22}\selectfont\mvSub};
  \node[anchor=north west, inner sep=0pt, text=white] at (24.5,141) {\sffamily\large\mvWhat};
  \node[anchor=north west, inner sep=0pt, text=mvbluemid] at (24.5,127) {\sffamily\mvVer};
}
\design{Mosaic of the screenshots, grayscale}{
  \fill[white] (0,0) rectangle (210,297);
  \mosaic
  \fill[white, blend mode=saturation] (0,0) rectangle (210,297);
  \fill[white, opacity=0.94] (0,104) rectangle (210,196);
  \draw[black, line width=0.6pt] (0,104) -- (210,104) (0,196) -- (210,196);
  \node[anchor=north west, inner sep=0pt, text=black] at (24,182) {\mvhead\fontsize{42}{46}\selectfont\mvTitle};
  \node[anchor=north west, inner sep=0pt, text=black!70] at (24.5,162) {\rmfamily\itshape\fontsize{17}{22}\selectfont\mvSub};
  \node[anchor=north west, inner sep=0pt, text=black] at (24.5,141) {\sffamily\large\mvWhat};
  \node[anchor=north west, inner sep=0pt, text=black!60] at (24.5,127) {\sffamily\mvVer};
}

% ------------------------------------------------------------------ 12 name dot
\design{Name dot: the title as a watch's label}{
  % a curve with its control points, drawn as V°Watch draws a watch
  \draw[watchmag, line width=1.6pt] (14,118) .. controls (60,190) and (110,60) .. (150,140) .. controls (165,170) and (185,150) .. (198,128);
  \foreach \p in {(14,118),(49,140.6),(150,140),(198,128)} \fill[watchmag] \p ++(-1.6,-1.6) rectangle ++(3.2,3.2);
  % the big dot, up and right of the point on the curve (49,140.6), its tail curving down from its bottom edge
  \fill[watchmag] (62,178) -- (80,178) .. controls (76,162) and (62,148) .. (49,140.6) .. controls (60,150) and (64,164) .. (62,178) -- cycle;
  \node[fill=watchmag, rounded corners=4mm, inner xsep=7mm, inner ysep=5mm, anchor=south west, text=white] at (56,176)
    {\mvhead\fontsize{40}{44}\selectfont\mvTitle};
  % a smaller one
  \fill[watchmag] (148,104) -- (156,104) .. controls (156,118) and (153,130) .. (150,140) .. controls (150,128) and (150,116) .. (150,104) -- cycle;
  \node[fill=watchmag, rounded corners=2mm, inner xsep=3mm, inner ysep=2mm, anchor=north west, text=white] at (138,106)
    {\sffamily\bfseries Manual \mvversion};
  \node[anchor=north west, inner sep=0pt, text width=160mm] at (24,74) {\mvlight\fontsize{18}{23}\selectfont\color{mvink}\mvSub};
  \node[anchor=north west, inner sep=0pt] at (24.5,56) {\sffamily\color{mvgray}\mvWhat};
  \node[anchor=south west, inner sep=0pt] at (24.5,20) {\sffamily\small\color{mvgray}\mvVer};
}

% ------------------------------------------------------------------ 13 Rhino viewport
\design{Rhino viewport: lesson 01's bow tie in a Top view}{
  \fill[rhinobg] (0,0) rectangle (210,297);
  \foreach \x in {0,5,...,210} \draw[rhinominor, line width=0.25pt] (\x,0) -- (\x,297);
  \foreach \y in {2,7,...,297} \draw[rhinominor, line width=0.25pt] (0,\y) -- (210,\y);
  \foreach \x in {0,25,...,210} \draw[rhinomajor, line width=0.45pt] (\x,0) -- (\x,297);
  \foreach \y in {2,27,...,297} \draw[rhinomajor, line width=0.45pt] (0,\y) -- (210,\y);
  \draw[red!75!black, line width=0.8pt] (0,102) -- (210,102);
  \draw[green!50!black, line width=0.8pt] (50,0) -- (50,297);
  \node[anchor=north west, inner sep=0pt] at (5,293) {\sffamily\small\color{black!75}Top};
  % the outline with two corners swapped: a bow tie
  \draw[watchmag, line width=1.4pt, line join=round] (50,102) -- (175,102) -- (50,180) -- (175,180) -- cycle;
  \foreach \p in {(50,102),(175,102),(50,180),(175,180)} \fill[watchmag] \p circle (1.6);
  \fill[watchmag] (70,198) -- (78,198) .. controls (78,190) and (60,186) .. (50,180) .. controls (60,188) and (70,192) .. (70,198) -- cycle;
  \node[fill=watchmag, rounded corners=1.5mm, inner xsep=2.5mm, inner ysep=1.2mm, anchor=south west, text=white] at (66,196.5) {\sffamily\small outline};
  \fill[watchmag] (195,120) -- (201,120) .. controls (201,112) and (185,106) .. (175,102) .. controls (186,108) and (195,113) .. (195,120) -- cycle;
  \node[fill=watchmag, rounded corners=1.5mm, inner xsep=2.5mm, inner ysep=1.2mm, anchor=south east, text=white] at (205,118.5) {\sffamily\small corners (4)};
  % the title, as a panel floating over the view
  \fill[black, opacity=0.12] (21,214) rectangle (195,282);
  \fill[white, opacity=0.94] (19,216) rectangle (193,284);
  \node[anchor=north west, inner sep=0pt] at (27,276) {\mvhead\fontsize{40}{44}\selectfont\color{mvblue}\mvTitle};
  \node[anchor=north west, inner sep=0pt] at (27.5,257) {\mvlight\fontsize{16}{20}\selectfont\color{mvink}\mvSub};
  \node[anchor=north west, inner sep=0pt] at (27.5,240) {\sffamily\color{mvgray}\mvWhat};
  \node[anchor=north west, inner sep=0pt] at (27.5,229) {\sffamily\small\color{mvgray}\mvVer};
  \node[anchor=north west, inner sep=0pt, text width=120mm] at (57,92) {\sffamily\small\color{black!70}Four corners that each look right
    in numbers. Drawn, the outline crosses itself: two of them are swapped.};
}

% ------------------------------------------------------------------ 14 code editor
\newcommand{\cl}[3]{% a code line: number, y, code
  \node[anchor=base east, inner sep=0pt] at (25,#2) {\ttfamily\small\color{vsgutter}#1};
  \node[anchor=base west, inner sep=0pt] at (31,#2) {\ttfamily\small\color{vstext}#3};}
\design{Code editor: the title as a paused script}{
  \fill[vsbg] (0,0) rectangle (210,297);
  \fill[white, opacity=0.04] (0,270) rectangle (210,284);
  \node[anchor=west, inner sep=0pt] at (8,277) {\sffamily\small\color{vstext}miniVivisection_Manual.py};
  \fill[vsfn!35!vsbg] (0,200.4) rectangle (210,206.2);
  \fill[red!80!black] (12,203.3) circle (1.6);
  % (\color swallows the space after it: spaces go before it)
  \cl{1}{258}{\color{vscom}\#! python 3}
  \cl{2}{250}{\color{vskw}import \color{vstext}Rhino.Geometry \color{vskw}as \color{vstext}rg}
  \cl{3}{242}{\color{vskw}import \color{vstext}miniVivisection \color{vskw}as \color{vstext}mv}
  \cl{4}{234}{}
  \cl{5}{226}{\color{vscom}\# See your script's data in Rhino while you debug}
  \cl{6}{218}{\color{vsvar}manual \color{vstext}= mv.\color{vsfn}Manual\color{vstext}(\color{vsvar}version\color{vstext}=\color{vsstr}"\mvversion"\color{vstext}, \color{vsvar}build\color{vstext}=\color{vsnum}\mvbuild\color{vstext})\quad \color{vsgutter}manual = Manual(…)}
  \cl{7}{210}{\color{vskw}for \color{vsvar}script \color{vskw}in \color{vsvar}manual\color{vstext}.scripts:\quad \color{vsgutter}script = 01_watch_geometry.py}
  \cl{8}{202}{\quad \color{vsvar}corners \color{vstext}= script.\color{vsfn}points\color{vstext}()\quad \color{vsgutter}corners = [(0,0,0),(8,0,0),(0,5,0),(8,5,0)]}
  \cl{9}{194}{\quad mv.\color{vsfn}watch\color{vstext}(\color{vsvar}corners\color{vstext})\quad\quad \color{vscom}\# drawn in Rhino, with its name}
  \cl{10}{186}{\quad mv.\color{vsfn}edit\color{vstext}(\color{vsvar}corners\color{vstext})\quad\quad\ \color{vscom}\# drag them, and the script goes on with yours}
  \node[anchor=north west, inner sep=0pt] at (31,150) {\ttfamily\fontsize{38}{42}\selectfont\color{vsfn}\mvTitle\color{vstext}()};
  \node[anchor=north west, inner sep=0pt] at (31,128) {\ttfamily\large\color{vscom}\# \mvSub};
  \node[anchor=north west, inner sep=0pt] at (31,118) {\ttfamily\large\color{vscom}\# \mvWhat};
  \fill[vsstatus] (0,10) rectangle (210,18);
  \node[anchor=west, inner sep=0pt] at (8,14) {\sffamily\small\color{white}miniVivisection \mvversion\ · build \mvbuild\ · \mvdate};
  \node[anchor=east, inner sep=0pt] at (202,14) {\sffamily\small\color{white}V°Debug ON \quad Paused at line 8};
}

% ------------------------------------------------------------------ 15 blueprint
\design{Blueprint: the mark as a technical drawing}{
  \fill[blueprint] (0,0) rectangle (210,297);
  \foreach \x in {0,5,...,210} \draw[white, opacity=0.07, line width=0.25pt] (\x,0) -- (\x,297);
  \foreach \y in {0,5,...,297} \draw[white, opacity=0.07, line width=0.25pt] (0,\y) -- (210,\y);
  \foreach \x in {0,25,...,210} \draw[white, opacity=0.16, line width=0.4pt] (\x,0) -- (\x,297);
  \foreach \y in {0,25,...,297} \draw[white, opacity=0.16, line width=0.4pt] (0,\y) -- (210,\y);
  \draw[white, opacity=0.7, line width=0.8pt] (8,8) rectangle (202,289);
  \node[inner sep=0pt, text=white] (v) at (102,190) {\vlogo\fontsize{300}{300}\selectfont\outlined{1}{V°}};
  % dimensions
  \begin{scope}[white, line width=0.35pt, font=\ttfamily\scriptsize]
    \draw ([yshift=-3mm]v.south west) -- ++(0,-14);
    \draw ([yshift=-3mm]v.south east) -- ++(0,-14);
    \draw[{Stealth[length=2mm]}-{Stealth[length=2mm]}] ([yshift=-14mm]v.south west) -- node[fill=blueprint, inner sep=1.5pt]{mini · Vivisection} ([yshift=-14mm]v.south east);
    \draw ([xshift=-3mm]v.north west) -- ++(-12,0);
    \draw ([xshift=-3mm]v.south west) -- ++(-12,0);
    \draw[{Stealth[length=2mm]}-{Stealth[length=2mm]}] ([xshift=-12mm]v.north west) -- node[fill=blueprint, inner sep=1.5pt, rotate=90]{V° = 300 pt} ([xshift=-12mm]v.south west);
    \draw[dashed] ([xshift=-23mm, yshift=-13mm]v.north east) circle (13);
    \draw ([xshift=-14mm, yshift=-4mm]v.north east) -- ++(16,14) -- ++(22,0) node[anchor=south east]{DETAIL A: THE MARK};
  \end{scope}
  % the title block
  \begin{scope}[white, line width=0.6pt]
    \draw (92,14) rectangle (196,64);
    \draw (92,48) -- (196,48) (92,30) -- (196,30) (150,14) -- (150,30);
    \node[anchor=north west, inner sep=0pt] at (95,62) {\ttfamily\tiny TITLE};
    \node[anchor=west, inner sep=0pt] at (95,54) {\mvhead\fontsize{22}{22}\selectfont\color{white}\mvTitle};
    \node[anchor=north west, inner sep=0pt] at (95,46) {\ttfamily\tiny SUBJECT};
    \node[anchor=north west, inner sep=0pt, text width=98mm] at (95,42) {\ttfamily\scriptsize\color{white}\mvSub.\\ \mvWhat.};
    \node[anchor=north west, inner sep=0pt] at (95,28) {\ttfamily\tiny VERSION};
    \node[anchor=west, inner sep=0pt] at (95,20) {\ttfamily\small\color{white}\mvversion\ (build \mvbuild)};
    \node[anchor=north west, inner sep=0pt] at (153,28) {\ttfamily\tiny DATE · SHEET};
    \node[anchor=west, inner sep=0pt] at (153,20) {\ttfamily\small\color{white}\mvdate\ · 1};
  \end{scope}
  \node[anchor=north west, inner sep=0pt] at (16,281) {\ttfamily\small\color{white}MANUAL · DRAWING No. MV-\mvbuild};
}

% ------------------------------------------------------------------ 16 anatomical plate
\design{Anatomical plate: the mark laid open}{
  \fill[paper] (0,0) rectangle (210,297);
  \draw[sepia, line width=1.4pt] (12,12) rectangle (198,285);
  \draw[sepia, line width=0.4pt] (14.5,14.5) rectangle (195.5,282.5);
  \node[anchor=north, inner sep=0pt] at (105,272) {\rmfamily\small\addfontfeature{LetterSpace=20}\color{sepia}PLATE I.};
  \node[anchor=north, inner sep=0pt] at (105,262) {\rmfamily\bfseries\fontsize{36}{40}\selectfont\color{sepia}\mvTitle};
  \node[anchor=north, inner sep=0pt] at (105,245) {\rmfamily\itshape\fontsize{14}{18}\selectfont\color{sepia}or, the Anatomy of a Paused Script};
  \draw[sepia, line width=0.4pt] (80,236) -- (130,236);
  % the specimen, with fine hatching behind it
  \begin{scope}
    \clip (105,152) circle (52);
    \foreach \i in {-110,-107,...,110} \draw[sepia, opacity=0.18, line width=0.25pt] (105+\i,100) -- (105+\i+60,210);
  \end{scope}
  \draw[sepia, line width=0.5pt] (105,152) circle (52);
  \node[inner sep=0pt, text=sepia] (v) at (103,150) {\vlogo\fontsize{190}{190}\selectfont V°};
  % the labels, as on an engraved plate
  \begin{scope}[sepia, line width=0.35pt, font=\rmfamily\itshape\small]
    \draw (76,182) -- (40,206) node[anchor=south east, inner sep=1pt]{a.};
    \draw (101,112) -- (52,96) node[anchor=north east, inner sep=1pt]{b.};
    \draw (141,186) -- (170,206) node[anchor=south west, inner sep=1pt]{c.};
    \draw (124,128) -- (172,100) node[anchor=north west, inner sep=1pt]{d.};
  \end{scope}
  \node[anchor=north west, inner sep=0pt, text width=150mm] at (30,82) {\rmfamily\small\color{sepia}
    \textit{a.} The left limb, where values are drawn in the viewports (V°Watch). \quad
    \textit{b.} The apex, where the script is held at a breakpoint and opened. \quad
    \textit{c.} The mark, a watch's name dot. \quad
    \textit{d.} The right limb, where geometry is edited and handed back (V°Edit).};
  \node[anchor=north, inner sep=0pt] at (105,52) {\rmfamily\itshape\small\color{sepia}Fig. 1. — The V° of miniVivisection, laid open. \mvSub.};
  \node[anchor=south, inner sep=0pt] at (105,22) {\rmfamily\small\addfontfeature{LetterSpace=8}\color{sepia}RHINO 8 · SCRIPT EDITOR · GRASSHOPPER · \mvversion\ · \mvdate};
}

% ------------------------------------------------------------------ 17-25: the outlined V° (design 6) in the typographic layout (design 3)
% The typographic layout's parts. \swissheader{ink}; \swisstitle{top of the title}{ink}{color of "Vivisection"}{subtitle color};
% \swissfooter{ink}{gray}.
\newcommand{\swissheader}[1]{%
  \node[anchor=south west, inner sep=0pt, text=#1] at (24,270) {\sffamily\small\addfontfeature{LetterSpace=18}\bfseries MANUAL};
  \node[anchor=south east, inner sep=0pt, text=#1] at (186,270) {\sffamily\small Version \mvversion};
  \draw[#1, line width=0.5pt] (24,266) -- (186,266);}
\newcommand{\swisstitle}[4]{%
  \node[anchor=north west, inner sep=0pt, text=#2] at (23,#1) {\fontsize{58}{60}\selectfont{\mvlight mini}{\color{#3}\mvhead Vivisection}};
  \node[anchor=north west, inner sep=0pt, text width=162mm, align=left, text=#4] at (24.5,#1-30)
    {\rmfamily\itshape\fontsize{20}{26}\selectfont See your script's data in Rhino\\while you debug};
  \draw[#2, line width=2.5pt] (24.5,#1-56) -- (54.5,#1-56);
  \node[anchor=north west, inner sep=0pt, text width=100mm, text=#2] at (24.5,#1-64) {\rmfamily\large A plug-in for the Script Editor\\
    and Grasshopper of Rhino~8};}
\newcommand{\swissfooter}[2]{%
  \draw[#1, line width=0.5pt] (24,58) -- (186,58);
  \node[anchor=north west, inner sep=0pt, text width=50mm, text=#1] at (24,53) {\sffamily\footnotesize\bfseries Rhino 8\\\mdseries\color{#2}Script Editor\\Grasshopper};
  \node[anchor=north west, inner sep=0pt, text width=50mm, text=#1] at (80,53) {\sffamily\footnotesize\bfseries Version \mvversion\\\mdseries\color{#2}build \mvbuild\\experimental};
  \node[anchor=north west, inner sep=0pt, text width=50mm, text=#1] at (136,53) {\sffamily\footnotesize\bfseries \mvdate\\\mdseries\color{#2}formerly\\ScriptEditorWatchViz};}
% The V° at #1 pt, outlined with a #2 bp stroke.
\newcommand{\vout}[2]{\vlogo\fontsize{#1}{#1}\selectfont\outlined{#2}{V°}}

\design{Fusion: the outlined V° in the typographic layout}{
  \swissheader{black}
  \swisstitle{252}{black}{black}{black}
  \node[anchor=south east, inner sep=0pt, text=black] at (188,64) {\vout{250}{1.2}};
  \swissfooter{black}{black!65}
}

\design{Fusion: the outlined V° behind the type, bleeding off the page}{
  \begin{scope}
    \clip (0,0) rectangle (210,297);
    \node[anchor=south east, inner sep=0pt, text=black!30] at (236,-30) {\vout{620}{0.8}};
  \end{scope}
  \swissheader{black}
  \swisstitle{252}{black}{black}{black}
  \swissfooter{black}{black!65}
}

\design{Fusion with color: a blue V, an orange °}{
  \swissheader{black}
  \swisstitle{252}{black}{mvblue}{black!80}
  \node[anchor=south east, inner sep=0pt] at (188,64)
    {\vlogo\fontsize{250}{250}\selectfont{\color{mvblue}\outlined{1.6}{V}}{\color{mvorange}\outlined{1.6}{°}}};
  \swissfooter{black}{black!65}
}

\design{Fusion: the V° on top, the type below}{
  \swissheader{black}
  \node[anchor=north, inner sep=0pt, text=black] at (105,258) {\vout{260}{1.2}};
  \swisstitle{150}{black}{black}{black}
  \swissfooter{black}{black!65}
}

\design{Fusion: the V° filled light gray and outlined}{
  \swissheader{black}
  \swisstitle{252}{black}{black}{black}
  \node[anchor=south east, inner sep=0pt, text=black!8] at (188,64) {\vlogo\fontsize{250}{250}\selectfont V°};
  \node[anchor=south east, inner sep=0pt, text=black] at (188,64) {\vout{250}{1.2}};
  \swissfooter{black}{black!65}
}

\design{Fusion on the grid: the V° set on the layout's columns and baselines}{
  \foreach \y in {58,63,...,266} \draw[black!7, line width=0.3pt] (24,\y) -- (186,\y);
  \foreach \x in {24,80,136,186} \draw[black!12, line width=0.3pt] (\x,14) -- (\x,282);
  \swissheader{black}
  \swisstitle{252}{black}{black}{black}
  \node[anchor=south west, inner sep=0pt, text=black] at (22,63) {\vout{250}{1.2}};
  \node[anchor=south east, inner sep=0pt, text=black!55] at (186,63) {\sffamily\tiny V° · Segoe UI Bold · 250 pt};
  \swissfooter{black}{black!65}
}

\design{Fusion, reversed: white on black}{
  \fill[black] (0,0) rectangle (210,297);
  \swissheader{white}
  \swisstitle{252}{white}{white}{white!80!black}
  \node[anchor=south east, inner sep=0pt, text=white] at (188,64) {\vout{250}{1.2}};
  \swissfooter{white}{white!60!black}
}

\design{Fusion with one color: only the ° filled orange}{
  \swissheader{black}
  \swisstitle{252}{black}{black}{black}
  \node[anchor=south east, inner sep=0pt] at (188,64)
    {\vlogo\fontsize{250}{250}\selectfont{\color{black}\outlined{1.2}{V}}{\color{mvorange}°}};
  \swissfooter{black}{black!65}
}

\design{Fusion: the V° in inline lettering, a white line inside a black one}{
  \swissheader{black}
  \swisstitle{252}{black}{black}{black}
  \node[anchor=south east, inner sep=0pt, text=black] at (188,64) {\vout{250}{4.2}};
  \node[anchor=south east, inner sep=0pt, text=white] at (188,64) {\vout{250}{1.8}};
  \swissfooter{black}{black!65}
}

\end{document}
